\section{强关联体系}
强管理材料是体系中电子-电子相互作用非常重要，因而不能用独立粒子近似(如\textrm{DFT})来近似地合理描述地体系。这些体系中往往包含有局域地部分填充电子的元素。而电子相关效应导致了金属-绝缘体转变，磁性和非常规超导等。为了模拟这一类体系，人们发展了一套超越\textrm{DFT}的计算理论和方法，如\textrm{DFT}+$U$，约束无规相近似\textrm{(constrained Random-Phase Approximation, cRPA)}和动力学平均场理论\textrm{(Dynamical Mean-Field Theory, DMFT)}

\subsection{约束无规相近似}
\subsubsection{能带结构}
通过计算\ch{NiO}的初基原胞计算基态波函数和未占据态，并在$12\times12\times12$的$\vec k$点网格上计算$p$和$d$电子态的投影\textrm{DOS}。

强管理材料是体系中电子-电子相互作用非常重要，因而不能用独立粒子近似(如\textrm{DFT})来近似地合理描述地体系。这些体系中往往包含有局域的部分填充$d$电子和$f$电子的元素。电子相关效应引起的现象包括金属-绝缘体转变，磁性和非常规超导。为了模拟这一类体系，人们发展了若干超越\textrm{DFT}的模型，如\textrm{cRPA}，\textrm{DFT+}$U$和\textrm{DMFT}。%第一部分涉及\textrm{cRPA}，第二部分涉及\textrm{DFT+}$U$和\textrm{DMFT}

\textrm{cRPA}是可以计算模型\textrm{Hamiltonian}的有效相互作参数$U$，$J$和$J^{\prime}$。\upcite{PRB77-085122_2008, PRB112-254102_2025}。其主要思想是忽略$GW$方法中特定目标电子态的屏蔽\textrm{Coulomb}相互作用$W$的屏蔽效应，\textcolor{red}{结果某些屏蔽\textrm{Coulomb}相互作用可以通过由模型\textrm{Hamiltonian}描述的局域基函数张开的目标空间来计算}。通常目标空间的维度较低(一般包括5个电子态波函数)，因此也可以适用于高精度的计算，如\textrm{DMFT}。

一个\textrm{cRPA}计算包含如下几步:
\begin{enumerate}
	\item \textrm{DFT}基态计算\textrm{(e01)}
	\item 未占据轨道的\textrm{Kohn-Sham}本征态计算\textrm{(e01)}
	\item \textrm{DOS}投影(包括胖能带)\textrm{(e01)}
	\item 本征态的\textrm{Wannier}函数投影\textrm{(e02中\ch{Ni}的$d$和\ch{O}的$p$轨道)}
	\item \textrm{cRPA}计算\textrm{(e03)}
\end{enumerate}
本算例中，首先将计算\ch{NiO}的基态，并产生用于能带分析的未占据本征态。

\ch{NiO}结构文件\textrm{POSCAR}

计算控制文件\textrm{INCAR}

采用默认的\textrm{PBE}泛函完成基态的\textrm{DFT}计算，因此\textrm{INCAR}参数就是标准的\textrm{DFT}计算:

\texttt{SYSTEM}是计算对象的名称；\texttt{NBANDS}指定计算中总的\textrm{Kohn-Sham}轨道的数目；\texttt{ISMEAR}=-1表示\textrm{Fermi}附近采用\textrm{Fermi}函数积分方案(后续\texttt{LFINITE\_TEMPERATURE}所必须的)；\texttt{SIGMA}是展宽参数；\texttt{EDIFF}是电子步自洽迭代收敛标准。

$\vec k$空间网格剖分方案是以$\Gamma$点为中心的$12\times12\times12$网格文件\textrm{KPOINTS}

对于\textrm{cRPA}计算，采用带\textrm{\_GW}名称版本的原子数据集。相比于传统的赝势，这样的原子数据集的势函数可以在更大的能量范围内符合原子散射性质。用户可以通过各元素的\textrm{POTCAR}中\texttt{TITEL}参数确定元素的原子数据集版本。

执行\textrm{Bash}命名，完成\textrm{VASP}计算。

计算结束后得到\textrm{DFT}基态，通过\textrm{Py4vasp}模块可以得到\textrm{HOMO}和\textrm{LUMO}的能量本征值差，并绘制\textrm{DOS}曲线。

\textrm{PBE}计算的带隙是\textrm{0.0505356060869957}

小的带隙表明\ch{NiO}很可能是金属，但是因为其中存在强相关的电子，它实际上是\textrm{Mott}绝缘体。

通过计算，得到产生未占据态轨道所必须的\textrm{WAVECAR}文件。

第二步，增加未占据的数目确定长波极限下的极化率(数据保存在\textrm{WAVEDER}文件中)。

计算控制参数文件\textrm{unoccupied/INCAR}

增加\texttt{NBANDS}的数目，完成\textrm{Hamiltonian}的单次精确对角化\textrm{\texttt{ALGO}=Exact}。增加\texttt{NBANDS}的数目，因此计算的能带不仅包括占据轨道，而且包括\textrm{cRPA}计算必须的未占据轨道。本次计算的\textrm{WAVECAR}和\textrm{WAVEDER}文件是实际\textrm{cRPA}是计算的输入文件。设置\texttt{LOPTICS}=\textrm{T}来计算频率相关的介电函数矩阵，设置\texttt{LFINITE\_TEMPERATURE}=\textrm{T}来计算光学矩阵元。

执行\textrm{Bash}命令，复制\textrm{DFT}计算的\textrm{WAVECAR}，完成\textrm{VASP}计算。

由于\texttt{NBANDS}增加，因此计算的不仅包括占据轨道也包括未占据轨道。未占据轨道是后续\textrm{cRPA}计算所必须的，首先选定相关的能带来构建局域的\textrm{Wannier}函数基组。

第三步计算\ch{NiO}的\textrm{DOS}并确定$p$-和$d$-轨道为主的能带。

计算控制文件\textrm{fatband/INCAR}

\texttt{EMIN}和\texttt{EMAX}定义了绘制\textrm{pDOS}的能量下限与上限，而\texttt{NEDOS}确定绘制\textrm{DOS}曲线的横轴网格点数。

执行\textrm{Bash}命令，复制上一步计算的\textrm{WAVECAR}和\textrm{CHGCAR}文件，完成\textrm{VASP}计算。

检查哪些能带分别具有\ch{NiO}中\ch{Ni}的$d$轨道和\ch{O}的$p$轨道的特征。读者将会在\textrm{Fermi}面附近8个能带具有上述轨道的特征，这些能带的编号依次为2-9，覆盖的能量范围是-8到\textrm{2eV}的范围。

从这部分电子态中，可以确定目标空间，选择上述-8到\textrm{2eV}的8个能带，混合了$s$-，$p$-和$d$-轨道的特征(在文件\textrm{PROCAR}中，包含$spd$-和每个轨道的位点的角动量投影波函数)。也可以只选局域在\ch{Ni}原子的5个$d$-电子态，不过这样构建，可利用的局域基函数也会更少。相反，当\ch{O}的$p$-电子态(其在\ch{O}原子周围离域)参与混合后，得到的\textrm{Wannier}函数能够更好地描述\textrm{Fermi}能附近的轨道。当然，也可在构建\textrm{Wannier}时包含$s$-电子态，不过$s$-电子态与近邻的$p$-电子态纠缠(位于$\sim\mathrm{eV}$附近)，因此如果选择了$s$-态，就很难保证不带入$p$-电子态。而且即便包含了$s$-电子态，其对\textrm{Fermi}能\textrm{(0eV)}附近的5个能带的贡献极小，基本可以忽略。

应用\textrm{Py4vasp}模块可以考察$p$-和$d$-电子态对\textrm{DOS}贡献的情况。

通过\textrm{Python}脚本中改变\textsc{NiO\_band.dos.plot(``d'')}或\textsc{NiO\_band.dos.plot(``p'')}变成\textsc{NiO\_band.dos.plot(``Ni'')}或\textsc{NiO\_band.dos.plot(``O'')}来检查\textrm{DOS}中哪些位置是局域电子态。

在\textrm{DOS}中，在\textrm{Fermi}能级\textrm{(0eV)}附近的某些位置的贡献由$p$电子态渐变为$d$电子态为主，在\textrm{-5eV}附近则是由$d$电子态渐渐为$p$电子态为主，说明选择的基函数中包含$p$电子态的重要性。此外，在$\sim-7\mathrm{eV}$附近，主要也是来自$d$-电子态的贡献，而在\textrm{0eV}附近则基本没有$s$-电子态的贡献，因此无须包含$s$-电子态。由8个\textrm{Blochl}能带可以构建8个\textrm{Wannier}轨道。%下一步将会学习生成用于\textrm{cRPA}计算的\textrm{Wannier}基函数。

\subsubsection{\rm{Wannier}投影}
对\ch{NiO}通过$12\times12\times12$$\vec k$空间网格，通过\textrm{Wannier90}选项产生\textrm{Wannier}函数。

\textrm{Wannier}函数是构建小的局域基组的重要工具\upcite{PRB56-12847_1997,PRB75-195121_2007}。\textrm{Blochl}轨道$\Psi_{n\vec k}^{\sigma}(\vec r)$与平移的周期函数$\mu_{n\vec k}^{\sigma}(r)$的关系是
\begin{displaymath}
	\Psi_{n\vec k}^{\sigma}(\vec r)=\mathrm{e}^{\mathrm{i}\vec k\cdot\vec r}\mu_{n\vec k}^{\sigma}(r)
\end{displaymath}
这里$n$是能带指标，$\vec k$是晶体动量，$\sigma$是自旋。而\textrm{Blochl}轨道的线性组合可以得到\textrm{Wannier}函数$|w_{i\vec R}^{\sigma}\rangle$
\begin{displaymath}
	|w_{i\vec R}^{\sigma}\rangle=\dfrac1{N_k}\sum_{n\vec k}\mathrm{e}^{-\mathrm{i}\vec k\cdot\vec R}T_{in}^{\sigma(\vec k)}|\Psi_{n\vec k}^{\sigma}\rangle
\end{displaymath}
这里$N_k$是$\vec k$点数目，而$T_{in}^{\sigma(\vec k)}$是平移矩阵。这里使用$T$也因为\textrm{DFT+}$U$计算中使用了$U$(即一直使用的$U$-标志:$U_{nm\vec k}$)

通常这样构造的基函数是局域的，而如果周期性影像间的相互作用则可以忽略，因此可以由位于$\vec R=0$的元胞中的\textrm{Wannier}函数来计算
\begin{displaymath}
	|w_{i\vec R}^{\sigma}\rangle=\dfrac1{N_k}\sum_{n\vec k}T_{in}^{\sigma(\vec k)}|\Psi_{n\vec k}^{\sigma}\rangle
\end{displaymath}
由类似\textrm{cRPA}的模型\textrm{Hamiltonian}，可以使用少量\textrm{Bloch}波函数构建小的子空间，而目标电子态主要集中在以化学势(或\textrm{Fermi}能)为中心的能级附近，而且空间上局域在某些原子/离子附近。模型\textrm{Hamiltonian}可以通过\textrm{Wannier}基函数表示目标电子态，方便地求解。

%本算例中通过\textrm{Wannier90}投影出\textrm{Wannier}轨道，并绘制出波函数并可视化其各部分。

\ch{NiO}结构文件\textrm{POSCAR}

计算控制参数文件\textrm{INCAR}

计算控制参数\texttt{LWRITE\_WANPROJ}确定是否将\textrm{Wannier}投影函数写入文件\textrm{WANPROJ}，该文件中包含每个$\vec k$点上指定投影的能带的\textrm{Wannier}投影矩阵\textbf{T}。\texttt{NUM\_WANN}控制须构建的\textrm{Wannier}函数的数目，\texttt{WANNIER90\_WIN}指定\textrm{wannier90.win}文件的内容，\textrm{LWANNIER90\_RUN}让\textrm{VASP}调用\textrm{wannier90}库函数。注意，由于计算控制参数\texttt{exclude\_bands}的规定，\textrm{Wannier}投影函数只包括8个能带。此外的\textrm{Wannier90}相关的计算控制参数选项确定了\textrm{Wannier}函数的数目，用于投影的能带数目和被投影的能带中的轨道类型。最后通过\textrm{gnuplot}文件绘制对应的胖能带图。设置计算控制参数\texttt{ALGO}=\textrm{NONE}，确保波函数不变，因为这是后处理步骤，波函数也不应该再有变化。

$\vec k$空间网格剖分方案是以$\Gamma$点为中心的$12\times12\times12$网格文件\textrm{KPOINTS}

执行\textrm{Bash}命令，复制上一步计算的\textrm{WAVECAR}文件，完成\textrm{VASP}计算。

由此会依据前面产生的8个\textrm{Bloch}轨道，产生8\textrm{Wannier}轨道(数据写入\textrm{WANPROJ})。在\textrm{INCAR}文件中，通过计算控制参数\texttt{NUM\_WANN}=8来确定8个电子态，\textcolor{red}{对于\textrm{Wannier90}处理的计算控制参数，包括：}

这些参数指定了参与\textrm{Wannier}投影的只限于\ch{Ni}的5$d$轨道和\ch{O}的3$p$为主的8个能带。

应用\textrm{gnuplot}可以根据数据绘制源于$d$和$p$电子态的分析图文件。

然后采用\textrm{Python}命令得到能带组分的分析图。

如果选择了\textrm{Wannier}电子态，采用\textrm{Wannier}基组得到的能带结构，将和前面采用\textrm{Blochl}基组的结果类似。

上述能带结构为\textrm{cRPA}计算很重要的内容。能带图中颜色的搭配反映了每个能带的原子轨道的特点。图中\ch{Ni}的$d$电子态标注为\textcolor{red}{红色}，而\ch{O}的$p$电子态标注为\textcolor{blue}{蓝色}。而\textrm{Fermi}能级附近($\sim0\mathrm{eV}$)的5个\textrm{Wannier}能带主要是$d$-电子特征，也有明显的$p$-电子态的贡献，而在更低能量($\sim-7\mathrm{eV}$)附近的\textrm{Wannier}能带，最重要的贡献来自$p$电子态，但也有明显的贡献来自$d$-电子态。

作为\textrm{cRPA}计算的对象，本算例选择的是\textrm{Fermi}能附近的$5d$-电子态，这里这些$d$电子态通过模型\textrm{Hamiltonian}描述，在\textrm{wannmerf90.wout}文件中看到\textrm{Wannier90}的输出结果。

从计算数据来看，前五个\textrm{Wannier}函数局域在位于\textsc{(0.0,0.0,0.0)}位置的\ch{Ni}原子附近，后三个\textrm{Wannier}函数局域在位于\textsc{(2.0855,2.0855,2.0855)}的\ch{O}原子附近。这些电子态对应与5$d$轨道和3$p$轨道。\textcolor{red}{如果想要选定最局域的轨道，即位于$\sim0.45$的\ch{Ni}为中心的$d$轨道和位于$\sim0.87$的\ch{O}为中心的$p$轨道。}%关于能带展宽，参见文献\cite{PRB56-12847_1997}。

如果对于仅含有$d$-电子轨道的效果有兴趣，可以单独完成\textsc{d-only}目录下的计算。

产生一个不同的\textrm{WANPROJ}文件，后续应用\textrm{cRPA}计算

执行\textrm{Bash}命令，复制考虑未占据体系的计算的\textrm{WAVECAR}文件，完成\textrm{VASP}计算。

\textcolor{red}{在\textrm{VASP}计算过程中，对照两种不同的计算控制参数文件\textrm{INCAR}}

构建的\textrm{Wannier}电子态数目由8降为5，不再包括3$p$能带，(即能带2-4，\texttt{exclude\_bands=1-4})，没有包括\ch{O}的$p$电子态的投影(\textsc{O:p}已去除)。

计算结束后，同样应用\textrm{Bash}命令在\textrm{wannier90.wout}文件中检查\textrm{Wannier}轨道的裂分扩展情况。

结果显示，$d$-轨道裂分的$t_{2g}(0.95788)$和$e_g(1.6914)$扩展有明显的区别。两者都比包含$p$-轨道特征时分别有了大2-4倍的扩展\textcolor{red}{(见上面:0.45)}。这也表明了为了生成\textrm{Wannier}基函数，包含$p$-轨道的必要性和重要性。

\subsubsection{约束无规相近似}
针对\ch{NiO}初基原胞，采用\textrm{cRPA}针对8个\textrm{Wannier}函数计算\textrm{Hubbard}$U$，$U^{\prime}$和$J$参数，对比$4\times4\times4$和$6\times6\times6$两套$\vec k$空间网格剖分方案和两类频率采样点条件下对\textrm{Coulomb}势函数的影响。

通过仔细选择\textrm{Wannier}轨道$w_a^{\sigma}(\vec r)$，可以计算\textrm{cRPA}相互作用。通过保存在\textrm{WANPROJ}文件中的\textrm{Wannier}基函数，可以给出有效相互作用矩阵$U_{ijkl}$。
\begin{displaymath}
	U_{ijklj}^{\sigma\sigma^{\prime}}=\int\mathrm{d}\vec r\mathrm{d}\vec r^{\prime}w_i^{\ast\sigma}(\vec r)w_j^{\sigma}(\vec r)U(\vec r,\vec r^{\prime},\omega)w_k^{\ast\sigma^{\prime}}(\vec r^{\prime})w_l^{\sigma^{\prime}}(\vec r^{\prime})
\end{displaymath}
这里$U(\vec r,\vec r^{\prime},\omega)$是频率$\omega$空间的\textrm{Coulomb}相互作用势。

实际上，\textrm{Hubbard-Kanamori}相互作用可以写成
\begin{displaymath}
	\begin{aligned}
		U^{\sigma\sigma^{\prime}}=&\dfrac1N\sum_{i\in T}U_{iiii}^{\sigma\sigma^{\prime}}\\
		U^{\prime\sigma\sigma^{\prime}}=&\dfrac1{N(N+1)}\sum_{i,j\in T,i\neq j}U_{ijji}^{\sigma\sigma^{\prime}}\\
		J^{\sigma\sigma^{\prime}}=&\dfrac1{N(N-1)}\sum_{i,j\in T,i\neq j}U_{ijij}^{\sigma\sigma^{\prime}}
	\end{aligned}
\end{displaymath}
该表达式中的求和限于单个晶胞(即$\vec R=0$; $V_{ijkl}$，$U_{ijkl}$)。对于扩展体系(即$R>0$)，对应的是\textrm{Cooulomb}相互作用的积分。

对于实际的\textrm{DFT+}$U$计算\upcite{PRB44-943_1991, PRB57-1505_1998}，采用的是一种特殊的平均方式，将在\textrm{DFT_DMFT}中讨论。

计算\ch{NiO}的结构文件\textrm{POSCAR}

计算控制参数文件分别在\textrm{./INCAR}文件和\textrm{fine\_grid/INCAR}文件中

为了完成\textrm{cRPA}计算，设置计算控制参数\texttt{ALGO}=\textrm{CRPA}。该计算要使用局域基，这一点可通过设置\texttt{LOCALIZED\_BASIS}来指定。如果计算目录下存在\textrm{WANPROJ}文件，\textrm{VASP}会使用该文件中的数据作为对应的局域基。在本算例中，使用前面计算得到的\textrm{Wannier}投影函数。\texttt{NTARGET\_STATES}确定\textrm{cRPA}计算中不考虑的\textrm{Wannier}态波函数。\textrm{LFINITE\_TEMPERATURE}将会在\textrm{cRPA}计算中考虑有限温度的影响(要求采用\textrm{Fermi}展宽，\texttt{ISMEAR}=-1)，这对于零带隙或纠缠体系是必要的。另外不要忘记的是，\texttt{NBANDS}必须与产生的\textrm{WAVECAR}文件时设置的值相同。

等离震荡计算的计算控制参数\textrm{plasmons/INCAR}

$\vec k$空间网格剖分形式，在\textrm{./KPOINTS}文件和\textrm{plasmons/KPOINTS}文件中，都采用以$\Gamma$点为中心，$4\times4\times4$的粗网格方案

密网格布点，在\textrm{fine\_grid/KPOINTS}文件中，是以$\Gamma$点为中心，$6\times6\times6$的方案

在粗网格布点目录下，执行\textrm{Bash}命令，复制包含未占据态的\textrm{WAVECAR}文件和\textrm{Wannier}投影计算中的\textrm{WANPROJ}文件，完成\textrm{VASP}计算。

在\textrm{OUTCAR}文件中，查找\textsc{Hubbard U}字符串。

注意，更详尽的数据在\textrm{vasprun.h5}文件中。

应用\textrm{Py4vasp}模块，分析各类相互作用$U$,$U^{\prime}$，$J$参数。

增加$\vec k$空间网格点数目，考察网格点数目对计算的影响。在密网格目录下，$6\times6\timrs6$网格比$4\times4\times4$网格更密。(注意，这里计算的网格必须是能被\textrm{DFT}计算使用网格数整除，如之前使用的是$12\times12\timrs12$，这里使用$8\times8\timrs8$的网格就不可以。)

类似地，在密网格目录下，执行\textrm{Bash}命令，复制包含未占据态的\textrm{WAVECAR}和\textrm{WAVEDER}文件，以及\textrm{Wannier}投影计算中的\textrm{WANPROJ}文件，完成\textrm{VASP}计算。

一旦计算结束，同样应用\textrm{Py4vasp}模块，分析各类相互作用$U$,$U^{\prime}$，$J$参数。

通过增加$\vec k$空间网格剖分点，净的\textrm{Coulomb}相互作用几乎没有变化，表明该相互作用随$\vec k$点网格已经收敛。后续的\textrm{DMFT}计算中，需要屏蔽\textrm{Hubbard}相互作用$U$，$U^{\prime}$和$J$的数值，所以这里要重点注意这些数据。类似地可以看到，计算中$J$也已经收敛，几乎不随着$\vec k$点增加发生变化。而$U$和$U^{\prime}$的变化(差不多有\textrm{0.1 eV})，看起来还没有收敛，所以后续计算采用密网格。

在\textrm{cRPA}计算中设置\texttt{LTWO\_CENTER}，该参数指定非单中心的\textrm{Coulomb}积分，因此在文件\textrm{vaspout.h5}文件中数据集\textsc{uijkl}包含\textrm{Wannier}函数中心\textsc{R}之外的维度。

该数据集存储如下积分(表达式中的$R$是描述\textrm{Coulomb}势在两个元胞之间的衰减形式)
\begin{displaymath}
	U_{ijkl}(R)=\int\mathrm{d}r\int\mathrm{d}r^{\prime}w_i(\vec r-\vec R)w_j(\vec r-\vec R)U(\vec r,\vec r^{\prime},\omega=0)w_k(\vec r^{\prime})w_l(\vec r^{\prime})
\end{displaymath}

\textrm{Hubbard-Kanamori}$U$函数的空间衰减通常采用\textrm{Ohno}势函数研究，即
\begin{displaymath}
	\dfrac{U(R)}{U(0)}=\sqrt{\dfrac{\delta}{R+\delta}}
\end{displaymath}
应用\textrm{Bash}命令查找\textrm{vaspout.h5}文件中的数据结构

通过\textrm{Py4vasp}模块可以绘制模拟\textrm{Coulomb}势函数曲线。调节$delta$达到给定数据。尝试改变\textsc{radius\_max}参数，使得曲线与各离散点尽可能一致。

图中每个点是晶体中\ch{Ni}原子的最近邻相互作用，但完全没有考虑\ch{O}原子，因为这里的\textrm{cRPA}计算仅考虑了5$d$轨道的贡献。

\textrm{Ohno}势主要是由\textrm{Hubbard}$U$势导出的，但\textrm{Ohno}势被长程的噪声数据点严重地扭曲。\upcite{TCA2-219_1964, NPJCM10-286_2024}。这里的噪声主要源自采样的$\vec k$点不够，因此噪声可以通过增加采样$\vec k$点得到压制。此外，拟合中扣除噪声数据，也会提升两者的一致性。按照新的步骤重新计算
\begin{enumerate}
	\item 由上述绘图曲线中提取数据
	\item 扣除异常值(比如上图中超过10\AA处的\textrm{0.74eV})
	\item 绘制滤除数据并拟合\textrm{Ohno}势函数
\end{enumerate}

优化的\textrm{Ohno}参数为0.09749

注意，在本算例中，为了说明拟合曲线如何完善，移除的点都是源自噪声。实际应用中，应该增加$\vec k$空间网格剖分点和平面波截断能，因为过于稀疏的\textrm{FFT}网格会产生噪声数据。

对于大的晶胞，可以通过\textrm{Euclidean}空间-时间算法\upcite{CPC117-21100174-X_1999, PRB94-165109_2016}来评估\textrm{cRPA}计算。该方法在虚时间轴上执行所有内部计算。从这个角度说，一般的\textrm{Minkowski}时空度规$\eta=(-+-+)$就变换成$\eta=(++++)$，因此变成\textrm{Euclidean}空间$R^4$，有效\textrm{Coulomb}内核由此得以稀疏表示。

相比于传统的在\textrm{Minkowski}时空域内的\textrm{cRPA}计算，该方法可以让用户使用更多内存研究大体系有效相互作用参数。

相比于\textrm{Minkowski}算法，这种计算非常吃内存，而且对小体系计算更耗时，但对于大体系计算则快得多。使用4阶\textrm{MPI}并行和$4\times4\times4$网格点的计算，要大约\textrm{8GB}的\textrm{RAM}。

注意，在计算过程中，屏幕输出\textsc{estimated memory requirement per rank ****.* MB, per node ****.* MB}后，会有较长时间的没有响应，也没有输出。但此次程序仍然在执行。

类似地，在等离态计算目录下，执行\textrm{Bash}命令，复制包含未占据态的\textrm{WAVECAR}和\textrm{WAVEDER}文件，以及\textrm{Wannier}投影计算中的\textrm{WANPROJ}文件，完成\textrm{VASP}计算。

对\textrm{vaspout.h5}文件检查，看看其保存了哪些结果。

计算确定了位于\texttt{NOMEGA}虚频点的\textrm{Hubbard}相互作用，并将结果保存在\textsc{uijkl}数据集中。这些数据在实际使用中几乎没有用处，真正感兴趣的结果是作为实频函数的\textrm{Hubbard}相互作用，实频依赖的\textrm{Hubbard}相互作用可以通过拟合虚频数据形成\textrm{Pad\'e}表达式的解析延拓得到。\textrm{Pad\'e}表达式形式为
\begin{displaymath}
	U(z)=\dfrac{(z-u_1)(z-u_2)\cdots(z-u_N)}{(z-v_1)(z-v_2)\cdots(z-v_M)} \quad u_1,v_1,\cdots,u_N,v_M\in C
\end{displaymath}
应用\textrm{Py4vasp}模块，采用\textrm{AAA}算法可分别确定表达式分子和分母中多项式的根$u_1,v_1,\cdots,u_N,v_M$\upcite{SC20-A1494_2018}。不过在分子和分母上可能会有相同的根，在实际计算中两式则相互抵消。这种现象就是\textrm{Froissart}二重态。因为数值噪声，因为这些值留存在表达式中，破坏了解析延拓。可以使用参数\textsc{clean\_up\_tol}清除这些二重态。

采用\textrm{Py4vasp}模块绘制\textrm{Coulomb}相互作用曲线。

图中的屏蔽相互作用$U$存在存在多个尖峰和起伏的结构。第一个峰(位于约$\mathrm{5,5eV}$处)差不多对应于\textrm{Fermi}能级附近的$p$-和$d$-电子态质心的距离(参照\textrm{DOS}和\textrm{Wannier}函数)，在较低能量激发饰，这是最容易出现的激发。即使改变计算控制参数\texttt{clean\_up\_tol}参数并未显著改变该峰的位置，表明这是一个真正的等离振荡峰，而并非源于\textrm{Froissart}二重态。在该峰出现之前\textrm{(<\textrm{5eV})}，有效\textrm{Coulomb}相互作用基本保持常数不变，说明对于\ch{NiO}来说，采用类\textrm{Hubbard}模型近似描述\textrm{Coulomb}相互作用势是一个合理的近似。

在高频部分，屏蔽\textrm{Coulomb}相互作用势$U$开始逼近净的\textrm{Coulomb}势，为了解析更高频率势函数的演化趋势，需要计算更多的采样点。

如果将\texttt{NOMEGA}的是指由12增大到24，可以看到该参数的增加对高频区势函数的数值解析的影响。%为了节约时间，这里并未完成该计算，读者可以自行尝试完成。

该计算大约费时\textrm{10min}作用，通过\textrm{Bash}命令复制相关文件并通过\textrm{sed}命令，修改\textrm{INCAR}中的文件，然后执行\textrm{VASP}计算。

计算结束后，应用\textrm{Python}脚本绘制势函数曲线图。

使用24个虚频点，\textrm{AAA}拟合在高频区域给出更多的势能曲线变化的信息。不难看出，在低频区，\textrm{50eV}以上位置，\texttt{clean\_up\_tol}参数的变化，使得势函数曲线的起伏程度表现得更显著，这表明\textrm{Froissart}二重态的真实存在。随着数据点的增加，数值噪音得到显著压制。鉴于极小极大等距法\textrm{(Minimax isometry method)}选定的虚频点已经足够好\upcite{PRB101-205145_2020}，所以这是增加\texttt{NOMEGA}的值并未显著提升计算结果。相反，增加$\vec k$点数目有助于\textrm{Coulomb}势曲线在高频区的稳定性(注意这里计算用的虚频采用点数为24而不是12)，可以将\texttt{clean\_up\_tol}设置成更小的值。%用4核\textrm{MPI}、12\textrm{GB}内存的配置，计算需要较长的时间。

可以忽略这些计算过程，重点关注结果分析，因为结果在对应的目录中。

增加虚频采样点使得两个等离振荡峰更加明显，在高频区的\textrm{Coulomb}相互作用势曲线也更稳定。两个等离振荡峰对应两个不同的跃迁。第一个峰位于\textrm{5.5eV}，基本对应$d$-电子态和$p$-电子态的中心的距离，而位于\textrm{50eV}附近的等离振荡的主峰对应的是$GW$近似(典型的如\textrm{cRPA})下的群体激发。在更高频率下，可以看到屏蔽\textrm{Coulomb}势$U$开始逼近净\textrm{Coulomb}势$V$。只有在高密的$\vec k$空间网格剖分点和更多虚频网格采样点条件下，可以更好地解析相互作用函数。

上述分析最重要的结论是解析连续不是一个黑盒，可以看得出，当改变$\vec k$点网格和频率采样点，第一等离振荡峰并未发生显著改变，说明该峰是可靠的。在模型\textrm{Hamiltonian}方法中，绝大部分关心的物理现象都发生在低频区。而高频区$(\sim50\mathrm{eV})$的峰是更难收敛。该峰基本对应于$\Gamma$-辐射，因此对我们没有太多的意义。更高频率更难收敛，所以结果是不可靠的。当用户完成此类计算时，不要纠结于这些。应该关注的重点放在低频区有清晰物理意义的峰上。

如果用户有时间，推荐使用仅有$d$电子特征的\textrm{Wannier}函数轨道重新完成上述计算。首先完成\textsc{d-only}目录下的计算，然后重新计算\textrm{WANPROJ}文件。

仅考虑$d$电子态时，在$4\times4\times4$$\vec k$点网格密度下，计算净的\textrm{Hubbard}相互作用势$U$值，由26.3724变为22.8494。而屏蔽\textrm{Hubbard}势$U$将由7.0806变为6.118。这意味着有一些相互作用遗失了。
\subsection{\rm{DFT+}$U$和\rm{DFT+DMFT}}
%这部分讲义按照当前库中\textrm{Workshop}的结构编次。
用\textrm{VASP}完成\textrm{DFT}和\textrm{DFT+}$U$计算，再用\textsc{solild\_dmft}软件\textrm{TRIQS}采用\textrm{VASP}投影的函数\textrm{(PLOs)}完成\textrm{DFT+DMFT}计算。关于\textsc{TRIQS}软件如何在本机上安装，见安装说明文件\textsc{install\_triqs.txt}文本。

\subsubsection{\rm{DFT}自洽计算和准备\rm{DMFT}准备}
对反铁磁态\textrm{(AFM)}晶胞\ch{NiO}完成非铁磁\textrm{DFT}计算；将局域轨道转换称投影局域轨道\textrm{(Projected Local orbitals, PLOs)}，产生动力学平均场理论\textrm{(DMFT)}计算所需要的输入文件；并从之前的\textrm{cRPA}计算中提取$U$和$J$的值。

\textrm{DFT+}$U$是在标准\textrm{DFT}基础上，不用增加很大计算量就可以较好地校正局域轨道(典型的$d$和$f$轨道)中强关联电子的唯相计算策略。标准的(半)局域交换-相关泛函无法很好地描述自相互作用，局域电子，以至于低估\textrm{Mott}绝缘体和某些导体的带隙。\textrm{DFT+}$U$通过引入类\textrm{Hubbard}的在位\textrm{Coulomb}修正$U$修正了局域轨道存在双电子占据轨道的此类问题。根据\textrm{Dudarev}等提出的旋转不变性表示\upcite{PRB57-1505_1998}，体系总能量可以表示为
\begin{displaymath}
	E=E_{\mathrm{DFT}}+\sum_I\left[\dfrac{U^I}2\sum_{m,\sigma\neq m^{\prime},\sigma^{\prime}}n_m^{I\sigma}n_{m^{\prime}}^{I\sigma^{\prime}}-\dfrac{U^I}2n^I(n^I-1)\right]
\end{displaymath}
这里$n_m^{I\sigma}$是位于$I$格点处轨道$m$，自旋为$\sigma$的占据电子数。第二项$\dfrac{U^I}2n^I(n^I-1)$是重复计算校正，这一项扣除了包含在$E_{\mathrm{DFT}}$计算中估算的相互作用能，防止该项被重复计算两次。

对于完全占据，方括号中的各项彼此抵消。因此在原子极限下，体系总能仍然保持$E_{\mathrm{DFT}}$不变，但是考虑到$U$值的贡献，对于分数占据的情况，体系变得更为局域，更符合真是占据的结果。虽然\textrm{DFT+}$U$简单且定性正确，但参数$U$的选择必须很审慎(如可以通过\textrm{cRPA}或线性响应)得到，但是该方法不能描述动态相关效应或\textrm{Kondo}效应的物理本质。所以必须采用\textrm{DFT+DMFT}才能描述这种情况。

本算例中，先对反铁磁\ch{NiO}的非自旋极化\textrm{DFT}计算，准备局域轨道，完成\textrm{DFT+}$U$计算，提取$U$值和$J$值。

反铁磁\ch{NiO}的结构文件\textrm{POSCAR}

计算控制文件\textrm{INCAR}中各参数的含义
\begin{itemize}
	\item \texttt{SYSTEM}:~计算体系的名字
	\item \texttt{NCORE}:~设置每个能带计算中，有多少个核参与\textrm{FFT}计算
	\item \texttt{KPAR}:~设置$\vec k$点并行，即并行计算中，一次处理多少个$\vec k$点
	\item \texttt{ISMEAR}:~设置展宽参数。在四面体方法及其\textrm{Bl\"ochl}校正时，不需要设置展宽参数
	\item \texttt{KSPACING}:~指定$\vec k$空间网格剖分方案，给的是$\vec k$点的距离
	\item \texttt{EDIFF}:~指定电子步收敛的条件。为了达到波函数和本征值的严格收敛条件，该值可以设得比较小(对应\textrm{OSZICAR}中残量\textsc{rms}的模)
	\item \texttt{NELMIN}:~指定完成电子步自洽迭代的最少步数
	\item \texttt{NBANDS}:~指定计算的\textrm{Kohn-Sham}轨道的数目
	\item \texttt{LORBIT}:~指定轨道在原子球上的投影方法
	\item \texttt{EMIN}:~指定\textrm{DOS}计算的能量下限
	\item \texttt{EMAX}:~指定\textrm{DOS}计算的能量上限
	\item \texttt{NEDOS}:~指定\textrm{DOS}计算的采样网格点
	\item \texttt{LMAXMIN}:~指定写入\textrm{CHGCAR}文件中的\textrm{PAW}基函数中参与电荷密度混合并计算电荷密度的最大角动量量子数$l$
	\item \texttt{LOCPROJ}指定原子轨道中进行投影的局域轨道
\end{itemize}

$\vec k$空间网格剖分方案按照\textrm{INCAR}文件中\texttt{KSPACING}参数指定方法实现

计算中对\ch{Ni}采用原子数据集是\textsc{Ni\_pv}，包括了与$d$电子有弱的杂化的重要的$p$价电子态。数据集的选择显著影响能带结构和晶格参数，所以选择正确的原子数据集非常重要。理论上，这里应该采用\textsc{Ni\_sv\_GW}原子数据集，但这并非是因为计算时间的缘故。

在计算目录下，采用\textrm{Bash}命令完成\textrm{VASP}计算。

该计算是非磁性的标准的\textrm{DFT}计算。

讨论\textrm{INCAR}文件中的重要计算控制参数。通过\textsc{LOCPROJ}参数指定对\ch{Ni}-$d$轨道的投影。此外，指定\texttt{LORBIT}=14，在能量\texttt{EMIN}到\texttt{EMAX}范围内优化投影函数。对\textrm{PAW}投影函数的投影，留给用户的选择权是，可以选择向哪个\textrm{PAW}投影子轨道投影，通过该选项可以在指定能量范围内自动优化投影轨道选择。计算中的能带数目略有增加，为了保证\textrm{Kohn-Sham}波函数的收敛，收敛标准也更严格。

在\textrm{OSZICAR}文件的最后，对于较小的晶胞，相关数据的形式为

\textrm{Py4vasp}模块绘制\ch{NiO}的结构文件

完成该标准\textrm{DFT}计算，在当前目录\textsc{e04\_scf}目录下，得到输出文件\textrm{vaspout.h5}和投影数据文件\textrm{LOCPROJ}和\textrm{PROJCAR}，由该计算可以构造\textrm{DMFT}计算所需要的输入文件\textrm{TRIQS/solid\_dmft}。

\subsubsection{\rm{PLO}的构建和转换}
在\textrm{DFT+DMFT}计算中，需要确定相关的子空间(本算例中是\ch{Ni}的$3d$轨道)。这里称为投影局域轨道\textrm{(PLOs)}。在\textrm{VASP}计算中，\textrm{PLOs}的构建方式与计算控制参数\texttt{LDAU}指定的将\textrm{Kohn-Sham}电子态波函数$|\Psi_{\nv\vec k}\rangle$投影到\textrm{PAW}提供的局域轨道上的方案相同:
\begin{displaymath}
	P_{L,\nv}^{\mathbf{R}}(\vec k)=\sum_i\langle\chi_L^{\mathbf{R}}|\phi_i\rangle\langle\tilde p_i|\tilde{\Psi}_{\nv\vec k}\rangle
\end{displaymath}
这里$|\chi_L^{\mathbf{R}\rangle}$是与关联的位置$\mathbf{R}$有关的局域基函数，$|\phi_i\rangle$是全电子波函数，$\tilde p_i$是标准\textrm{PAW}投影子函数。$L$是复合的局域量子数。最初的投影子函数是正交的，在$\vec k$空间点$\vec k$，通过局域轨道数$m$与\textrm{Kohn-Sham}基函数$\nv$相关联:~$P_{m,\nv}^{\mathbf{R}}(\vec k)$。

局域投影函数有计算控制参数\texttt{LOCPROJ}选定，写入文件\textrm{LOCPROJ}(这些数据将会转换称\textrm{PLOs})和\textrm{PROJCAR}文件(该文件中包括格式输出数据)。通过\textrm{TRIQS/dft\_tools}处理初始的投影函数。\textsc{plo.cfg}定义了哪一个离子的哪个轨道作为关联轨道，而\textrm{Kohn-Sham}的轨道/能量范围将会投影到这些关联轨道上。\textcolor{red}{\textrm{plo}转换说明，有详细说明。}\textrm{Python}脚本\textrm{converter.py}(在\textrm{TRIQS/dft\_tools})读取\textrm{VASP}计算输出文件，并构建\textrm{vasp.h5}(这是\textrm{DFT+DMFT}软件\textrm{TRIQS/solid\_dmft}的输入文件。)

通读\textrm{plo.cfg}代码，在代码中该文件完成的工作是\textrm{TRIQS/DMFT}建立局域投影子向\textrm{PLOs}的转换信息。在本算例的\textrm{plo.cfg}文件中，关键的计算控制参数是
\begin{itemize}
	\item \textsc{EWINDOW}:~找到与\textrm{Kohn-Sham}电子态相关的能量范围
	\item \textsc{BANDS}:~用于投影的\textrm{VASP}计算的能量范围内的能量下标，\textcolor{red}{\textrm{VASP}中原始投影子是对全部能带的}
	\item \textsc{LSHELL}=2:~表示$d$-壳层
	\item \textsc{IONS}:~指定哪些原子带有强关联的$d$-壳层(两个反铁磁的\ch{Ni})
	\item \textsc{CORR}=\textrm{TRUE}，确保该壳层是正确的(即，由\textrm{DMFT}要处理的轨道)
\end{itemize}

在当前计算目录下，%执行\textrm{Bash}命令，

执行\textrm{Python}脚本\textrm{converter.py}
